% !TeX root = ../main.tex
\chapter{Challenges and Solutions for Using NAS in FL}\label{chapter:nas_in_fl_ch_and_sol}


We identified six overarching themes of challenges in the NAS-in-FL literature: 

We identified four overarching themes of challenges in the literature related to FL system characteristics: \textit{Client Compute Capabilities}, \textit{Client Networking Capabilities}, \textit{Client Heterogeneity} and the \textit{Trust and Threat Model}. Additionally, we identified one theme of challenges where the dominant source of the challenge relates to how the architecture search is coordinated: \textit{Architecture Search Coordination}.

For these themes we identified x sub-themes representing related challenges. For each sub-theme we describe concrete challenge themes relating to that sub-theme and the recurring solution types for that challenge in the literature.

TODO: solutions are organised under primary challenge they solve
TOD: a solution represents a distinct mechanism, conceptually similar mechanism are merged => how is this done?

The severity and applicability of a challenge as well as proposed solutions can vary slightly depending on the search strategy and search space used by the NAS-in-FL method (TODO: performance evaluation is basically universally trainig for short rounds). When this is the case we qualify the affected challenges and solutions with the search strategy and search space for which they are relevant (TODO: too strong a statement, should be about what is in literature instead).

For each challenge, we describe the most generally applicable solutions to each challenge first.


- flow chart showing applicability and severity of challenges and solutions at the end of chapter? => maybe appendix
- FL system characteristics for which the challenges are applicable/more severe
- NAS types based on NAS in FL and NAS 1000 papers categorisation

- client compute constraints
- client networking constraints
- client heterogeneity
- trust and threat model
- architecture search coordination